\section{Even operator walks and closed-walk signs}
\label{sec:tpc68-operator-walks}

Pair and Perron certificates take absolute values before different row
pairs interact.  Powers of the signed matrix provide a later place to take
absolute values.

\subsection{Power-row certificate}

For an integer \(q\ge1\), define
\begin{equation}
 \Omega_{2q}(H)
 :=
 \max_m\sum_n|(H^{2q})_{mn}|.
 \label{eq:tpc68-Omega-power}
\end{equation}

\begin{theorem}[Even operator-walk certificate]
\label{thm:tpc68-even-walk-certificate}
For every finite Hermitian matrix \(H\) and every \(q\ge1\),
\begin{equation}
 \boxed{
 \|H\|_{2\to2}
 \le\Omega_{2q}(H)^{1/(2q)}.}
 \label{eq:tpc68-even-walk-bound}
\end{equation}
\end{theorem}

\begin{proof}
The matrix \(H^{2q}\) is Hermitian, and its maximum absolute row and column
sums agree.  Therefore
\[
 \|H\|^{2q}
 =\|H^{2q}\|_{2\to2}
 \le\sqrt{\|H^{2q}\|_1\|H^{2q}\|_\infty}
 =\Omega_{2q}(H).
\]
Taking \(2q\)-th roots proves the result.
\end{proof}

The case \(q=1\) already includes interference between two-edge walks:
\[
 (H^2)_{mn}=\sum_kH_{mk}H_{kn}.
\]
It is not identical to the star estimate.  In particular,
\((H^2)_{mm}=s_m(H)^2\), while off-diagonal entries may cancel.

\subsection{Even moments}

Let \(r=\dim S_M\).

\begin{theorem}[Even-moment certificate]
\label{thm:tpc68-even-moment}
For every \(q\ge1\),
\begin{equation}
 \boxed{
 \|H\|
 \le\bigl(\tr H^{2q}\bigr)^{1/(2q)}
 \le r^{1/(2q)}\|H\|.}
 \label{eq:tpc68-even-moment-bound}
\end{equation}
\end{theorem}

\begin{proof}
If \(\lambda_1,\ldots,\lambda_r\) are the real eigenvalues of \(H\), then
\[
 \tr H^{2q}=\sum_{\nu=1}^r|\lambda_\nu|^{2q}.
\]
This sum is at least \(\max_\nu|\lambda_\nu|^{2q}\) and at most \(r\)
times that maximum.
\end{proof}

For a fixed finite matrix, the middle quantity decreases to \(\|H\|\) as
\(q\to\infty\).  In an asymptotic family whose dimension grows with \(X\),
using \(q=q_X\to\infty\) requires estimates uniform in both \(X\) and
\(q_X\).  A fixed-moment estimate cannot silently be extrapolated to that
regime.

\subsection{Literal decorated-walk expansion}

For \(m\ne n\), abbreviate each literal summand in
\eqref{eq:tpc68-literal-pair-list} by
\[
 a_{mn}(\xi):=\varepsilon_\xi^{mn}v_\xi^{mn},
 \qquad \xi\in\mathcal C_{mn}.
\]
For a collision pair \(\xi=(w,w')\), define
\[
 j(\xi):=j(w)=j(w'),
\]
where equality follows from the common complete native key.
Set \(\mathcal C_{mm}=\varnothing\), consistently with \(H_{mm}=0\).

\begin{proposition}[Exact decorated-walk identity]
\label{prop:tpc68-decorated-walk}
For every integer \(k\ge1\),
\begin{equation}
 (H^k)_{mn}
 =
 \sum_{\substack{m=m_0,m_1,\ldots,m_k=n\\
                 \xi_\ell\in\mathcal C_{m_{\ell-1}m_\ell}}}
 \prod_{\ell=1}^k
 a_{m_{\ell-1}m_\ell}(\xi_\ell).
 \label{eq:tpc68-decorated-walk-expansion}
\end{equation}
Consequently
\begin{equation}
 \tr H^{2q}
 =
 \sum_{\substack{m_0=m_{2q}\\
                 m_0,m_1,\ldots,m_{2q-1}\\
                 \xi_\ell\in\mathcal C_{m_{\ell-1}m_\ell}}}
 \prod_{\ell=1}^{2q}
 \varepsilon_{\xi_\ell}^{m_{\ell-1}m_\ell}
 v_{\xi_\ell}^{m_{\ell-1}m_\ell}.
 \label{eq:tpc68-literal-closed-walk-moment}
\end{equation}
Both identities sum the complete literal native collision lists on every
edge.
\end{proposition}

\begin{proof}
Expand the matrix product
\[
 (H^k)_{mn}
 =\sum_{m_1,\ldots,m_{k-1}}
 H_{mm_1}H_{m_1m_2}\cdots H_{m_{k-1}n}
\]
and then expand every entry by
\eqref{eq:tpc68-literal-pair-list}.  Summing the diagonal at \(k=2q\)
gives \eqref{eq:tpc68-literal-closed-walk-moment}.
\end{proof}

\begin{proposition}[Row gauges telescope on closed walks]
\label{prop:tpc68-gauge-telescoping}
Let \(u_m\in\mathbb C\) with \(|u_m|=1\), let
\(U=\diag(u_m)\), and put \(\widehat H=U^*HU\).  In every closed-walk term
of \(\tr\widehat H^{2q}\), the product of the row-gauge factors is one.
Hence
\[
 \tr\widehat H^{2q}=\tr H^{2q}.
\]
For the real self-adjoint M\"obius row gauge \(U_d=U_d^*=U_d^{-1}\), the
same statement applies to \(U_dHU_d\).
\end{proposition}

\begin{proof}
An edge \(m_{\ell-1}\to m_\ell\) contributes
\(\overline{u_{m_{\ell-1}}}u_{m_\ell}\).  Their product around a closed
walk is
\[
 \prod_{\ell=1}^{2q}
 \overline{u_{m_{\ell-1}}}u_{m_\ell}=1
\]
because \(m_{2q}=m_0\).  The trace identity also follows from unitary
invariance.
\end{proof}

\begin{proposition}[Pure same-residual walks have no target-sign gain]
\label{prop:tpc68-pure-residual-walk-sign}
Consider one decorated closed-walk term in
\eqref{eq:tpc68-literal-closed-walk-moment}.  If every decorated edge
\((w_\ell,w'_\ell)\) is a same-residual collision, then
\begin{equation}
 \prod_{\ell=1}^{2q}
 \varepsilon_{\xi_\ell}^{m_{\ell-1}m_\ell}=1.
 \label{eq:tpc68-pure-residual-sign-product}
\end{equation}
Therefore pure \(H_{\rm res}\) walks have no cancellation from the target
M\"obius signs themselves.  Their literal complex weights may still
cancel, and walks containing \(H_{\rm erase}\) edges retain
unequal-residual sign information.
\end{proposition}

\begin{proof}
By \cref{prop:tpc68-same-residual-sign-freeze}, the edge sign is
\[
 \varepsilon_{\xi_\ell}^{m_{\ell-1}m_\ell}
 =\mu(d_{m_{\ell-1}})\mu(d_{m_\ell}).
\]
Every vertex gauge occurs twice in the closed product, so the product is
one.
\end{proof}

The telescoping proposition explains why the common residual/source gauge
is not the arithmetic cancellation sought here.  It disappears from every
closed walk automatically.  The signs
\[
 \prod_{\ell=1}^{2q}
 \mu(m_{\ell-1}j(\xi_\ell)+h_0)
 \mu(m_\ell j(\xi_\ell)+h_0)
\]
in \eqref{eq:tpc68-literal-closed-walk-moment} are the genuine target
signs.  The proposition above identifies the pure same-residual subfamily
on which they are trivial.  Bounding the remaining complete weighted sum
is an L2 problem.

\subsection{Why a support walk count is insufficient}

Taking absolute values term by term in
\eqref{eq:tpc68-literal-closed-walk-moment} reduces the method to a
positive walk count and loses its intended advantage.  Conversely, a
random-sign heuristic for the decorated walks is not a proof for the one
literal M\"obius array.  A valid operator-walk input must estimate the
complete fixed-\(h_0\) weighted sum, uniformly over all start rows needed
by \(\Omega_{2q}\), or the complete closed-walk sum needed by
\(\tr H^{2q}\).
